\subsection{An idempotent image of the nilpotent opposite actor}

\begin{proposition}[Proved]
\label{CP2:A-nilpotent-four-idempotent-action}
Let
\[
 N=\{1,e_0,e_1,a\},\qquad
 e_i e_j=e_j,\qquad ae_i=e_i,\qquad e_i a=e_0,\qquad a^2=e_0,
\]
and let $S$ be a finite nonempty left-$N$, right-$M$ biset.  Suppose
the left actor $F$ of $a$ is idempotent.  If $|M|\leq11$ and the
collage has bounded rational homology, all its rational degree-two
products vanish.  There is no restriction on $|S|$ or on the second
Betti number.  The opposite assertion holds as well.
\end{proposition}
\begin{proof}
Write $P_0,P_1$ for the actors of $e_0,e_1$.  The action relations
give $F^2=P_0$, so idempotence of $F$ is equivalent to $F=P_0$.
Put $T=P_0S=P_1S$.  This is a nonempty $M$-invariant subset, fixed
pointwise by all of $N$.

We use the complete natural right-projective resolution of the
trivial $\mathbb QN$ module displayed in
Proposition~\ref{CP2:A-nilpotent-four-corner-four-states}.  Its
resolution and augmented chain comparison have no state-count
restriction.  For $R=\mathbb QM$, the reduced augmented right-$R$
complex is
\[
 C_0=L_T:=\ker(\mathbb QT\to\mathbb Q),\qquad
 C_j=D:=(1-P_0)\mathbb QS\quad(j\geq1),
\]
\[
 d_1=P_1-P_0,\qquad d_j=F-P_0=0\quad(j\geq2).
\]
Since $BN$ is rationally acyclic, the actual collage has reduced
chains quasi-isomorphic to
$(C\otimes_R^{\mathbf L}\mathbb Q)[1]$.  The chain comparison
also respects the inclusion of the actual sub-biset $T$, whose
reduced augmented complex is $L_T$ in degree zero.

There is an actual direct-sum decomposition of right-$R$ complexes
\[
 C=[D\xrightarrow{P_1-P_0}L_T]
       \ \oplus\ \bigoplus_{j\geq2}D[j].
\]
Tensor with a bounded-below projective left-$R$ resolution of
$\mathbb Q$.  Direct sums and totalization commute here; in each
total degree only finitely many shifts contribute.  If
$\operatorname{Tor}^R_i(D,\mathbb Q)$ were nonzero for any $i\geq0$,
the displayed direct summands would produce nonzero homology in
arbitrarily high degrees.  Boundedness therefore forces
$D\otimes_R^{\mathbf L}\mathbb Q$ to be acyclic.  Consequently
the inclusion $L_T\to C$ is a quasi-isomorphism after derived
tensoring.

It follows that inclusion of the actual subcollage on $T$ is a
rational homology equivalence, and hence induces a rational
cohomology-ring isomorphism.  On that subcollage the left $N$
action is trivial.  The actual functor replacing the corner $N$
by a point has bridge map
$BN\times B(T,M)\to B(T,M)$; because $BN$ is rationally acyclic,
it is a rational homology and ring equivalence onto the ordinary
$M$-act cone on $T$.  That cone inherits bounded rational homology.
Theorem~\ref{CP2:A-bounded-eleven-acts} makes all its rational
degree-two products zero.  Naturality transfers the conclusion
back to the original collage.  Opposition proves the dual case.
\end{proof}

\begin{remark*}
In a literal action search, this gate is precisely the test
$F^2=F$ for the image of the unique nonidempotent element $a$ of
$N$.  Its equivalence with $F=P_0$ uses the already imposed table
equation $a^2=e_0$.  The predicate is invariant under relabelling
the states.  Removing it from a finite scope requires exact counts
and independently checked control subsequences; the theorem alone
does not certify an unfinished search.
\end{remark*}
